\chapter*{Preface}

Probability theory studies the mathematical framework for quantifying random events and their likelihoods, while stochastic processes extend this framework to model systems that evolve randomly over time. The core concepts of this field are crucial for understanding phenomena encountered in engineering, such as communication systems, signal processing, and machine learning.

This notebook includes the necessary definitions, theorems, and examples for the 3412110099 (or BBC4941) Probability Theory and Stochastic Processes module, which I am taking in the spring semester of 2026 as part of the Joint Programme between Beijing University of Posts and Telecommunications (BUPT) and Queen Mary University of London (QMUL). Most of the content is based on the lecture slides provided by Professor YANG Jiankui at BUPT. GenAI tools have been used to assist in writing some of the solutions and explanations.

This notebook comprises nine chapters and an index, structured as follows:
\begin{description}
  \item[Chapter 0] reviews some basic concepts in calculus and complex analysis, which are necessary for understanding the content in the following chapters.
  \item[Chapter 1] gives the definition of probability and some of its properties, and introduces the concept of conditional probability and independence, which are fundamental for understanding how events relate to each other.
  \item[Chapter 2] defines random variables, distributions, and their numerical characteristics. Some common discrete and continuous distributions are also introduced, which are widely used in modeling random phenomena.
  \item[Chapter 3] extends the concept of random variables to random vectors, and introduces the joint distribution, marginal distribution, and conditional distribution of random vectors. The covariance matrix and correlation coefficient are also defined, which are important for understanding the relationships between different random variables.
  \item[Chapter 4] presents some limit theorems, which are fundamental results in probability theory that describe the behavior of sums of random variables.
  \item[Chapter 5] introduces stochastic processes by defining them and discussing their properties and classifications.
  \item[Chapter 6] focuses on stationary stochastic processes, which are a special class of stochastic processes that have statistical properties that do not change over time. The power spectral density is also introduced.
  \item[Chapter 7] introduces discrete-time Markov chains, which are a type of stochastic process that has the Markov property, meaning that the future state depends only on the current state and not on the past states.
  \item[Chapter 8] introduces independent increment processes, along with the Poisson process, Gaussian process, Brownian motion, and Wiener process, which are important examples of stochastic processes with independent increments.
  \item[Index] serves as a quick reference for the key terms and theorems.
\end{description}
Please note that Chapter 4 and the following chapters are still under construction and review, and will be updated in the future.

I hope this notebook can be a useful resource for students taking this module, and I welcome any feedback or suggestions for improvement. Let's enjoy the journey of learning probability theory and stochastic processes together! (\^{}~--~\^{})/'

\vspace{2em}
\begin{flushright}
  Henry Tan\\
  \href{https://github.com/Henry0512}{https://github.com/Henry0512}\\
  \today
\end{flushright}